\begin{lemma}[First-entry clipping of exact incidence paths]
Let $P$ be a simple $(\alpha,\beta)$ incidence walk from vertex $a$ to
vertex $b$ in a finite loopless multigraph with partial edge coloring $\varphi$.
Let $U$ be an exact tree sequence, with $a$ outside $U$ and $b$ in $U$.
Suppose $a$ misses at least one of the two colors.  There exist a vertex
$v$ and an exact walk $Q$ from $a$ to $v$ such that $Q$ is a
$(\alpha,\beta)$ exit path of $U$, its support is contained in the
support of $P$, and $|Q|\leq |P|$.  Moreover, for every exact edge $f$
whose edge-node occurs in $P$, if every endpoint of $f$ belongs to $U$,
then $|Q|<|P|$.
\end{lemma}
\begin{proof}
Consider the finite list of nodes in $P$, in its given direction.
The list contains a vertex of $U$ because its last node is $b$.
Choose the first occurrence of a vertex-node in $U$, call its vertex
$v$, and take $Q$ to be the walk prefix ending there.  Its initial
vertex $a$ is outside $U$, so $a\ne v$.  The adjacency rule in
\noderef{TwoColorIncidenceGraph} only permits transitions between
vertex-nodes and edge-nodes.  A nonempty walk between distinct
vertex-nodes therefore has length at least two: length zero would
identify its endpoints, and length one would give forbidden adjacency
between two vertex-nodes.  A prefix of a simple path is simple.  The
choice of the first entry implies that the only vertex of $U$ appearing
in $Q$ is $v$.  Its initial missing-color alternative is the given one.
These facts are exactly the clauses of \noderef{TreeExitPath}.
The node list of a walk prefix is a prefix of the original list, hence
its support is contained in the original support and its length is no
larger.

Fix an edge $f$ with the stated properties.  Its edge-node cannot occur
in $Q$.  Indeed, if it occurred at position $j$, neither endpoint of
$Q$ is an edge-node, so $0<j<|Q|$.  Its predecessor is a vertex-node
$w$, incident with $f$ by \noderef{TwoColorIncidenceGraph}.  By the
endpoint hypothesis $w$ is in $U$, so the first-entry property forces
$w=v$.  But this is an occurrence of $v$ before the last position of
$Q$, contradicting simplicity.  Thus $f$ occurs in $P$ and not in
$Q$.  If the two lengths were equal, the walk prefix $Q$ would contain
every position of $P$, including that occurrence of $f$.  Therefore
equality of lengths is impossible; together with the weak length
inequality this proves the strict inequality.
\end{proof}
